\section{GCN-Guided Adaptive Large Neighborhood Search}
\label{sec:method}
\subsection{Search Architecture}
GCN-ALNS separates three roles: learned assignment preferences, inexpensive
search guidance, and full scenario verification. Figure~\ref{fig:overview}
shows this separation. The learned component ranks candidate placements; it
does not replace the recovery simulator or alter the SLA definition.

For a candidate deployment, the quick evaluator computes its nominal
schedule, deployment cost $J_0$, and an approximate recovery-delay surrogate
$\rhat$. The surrogate takes the largest positive difference between a
task-local recovery estimate and that task's schedule slack. For an active
task, the local estimate includes detection, switching, synchronization, and
the finish-time gap between its two copies. For a passive task, it includes
detection, replanning, a queue estimate, incoming communication, and
reexecution. This approximation does not enumerate~\eqref{eq:fault_set} and
is not a proven upper bound on $\rmax$.

Define estimated violation and penalized cost as
\begin{align}
 V(\chi)&=\pos{\cbase-D}+\pos{\rhat-\reff},\\
 \Phi(\chi)&=J_0(\chi)+\eta_D\pos{\cbase-D}
                     +\eta_R\pos{\rhat-\reff}.
 \label{eq:quick_objective}
\end{align}
Before a fully verified feasible solution has been found, the search
minimizes $V+\epsilon J_0$, with $\epsilon=10^{-4}$ providing a small cost
preference. It switches to $\Phi$ after establishing feasibility. This
mechanism distinguishes the search's constraint-reduction phase from its
cost-reduction phase.

Full verification is reserved for candidates that the surrogate regards as
feasible and whose quick deployment cost is below the stored feasible
solution's reported score. If a candidate passes exhaustive scenario
enumeration, the feasible-solution register is updated with its deployment
and full cost. The quick screen uses $J_0$, whereas the stored score includes
the recovery proxy in~\eqref{eq:objective}; the resulting heuristic is not a
monotone descent method for the full reported objective. Moreover, an
overconservative surrogate can screen out an actually feasible deployment.
We therefore claim neither search completeness nor an optimality guarantee.

\begin{figure*}[!t]
 \centering
 \resizebox{0.97\textwidth}{!}{%
\begin{tikzpicture}[
 box/.style={draw,rounded corners=1pt,align=center,text width=3.8cm,
             minimum height=1.05cm,inner sep=5pt},
 arr/.style={-{Stealth[length=2.2mm]},thick},
 font=\small]
 \node[box] (train) at (0,2.3) {Fixed train/validation inputs\\No solver-test demonstrations};
 \node[box] (teacher) at (5.1,2.3) {Teacher panel and full verifier\\Feasible placement labels};
 \node[box] (model) at (10.2,2.3) {Resource GCN and task MLP\\Train then freeze scorer};
 \node[box] (test) at (0,0) {Held-out workflow and resources\\Independent evaluation input};
 \node[box] (search) at (5.1,0) {GCN-guided ALNS\\Quick score, SA, restarts};
 \node[box] (verify) at (10.2,0) {Full scenario verification\\Verified feasible solution register};
 \draw[arr] (train)--(teacher);
 \draw[arr] (teacher)--(model);
 \draw[arr] (test)--(search);
 \draw[arr] (search)--node[above,font=\footnotesize]{promising candidates}(verify);
 \draw[arr,dashed] (model.south) to[out=-90,in=90]
    node[above left,font=\footnotesize,align=center]{cached pair\\preferences}(search.north);
 \draw[arr] (search.south) to[out=-115,in=-65,looseness=3]
    node[below,font=\footnotesize]{destroy / repair loop}(search.south);
\end{tikzpicture}%
}

 \caption{Training and deployment workflow. Only fixed train/validation
 instances provide teacher demonstrations. The frozen scorer supplies
 preferences to ALNS on a held-out instance. A quick evaluator guides search,
 while full scenario enumeration determines whether a feasible solution may
 be retained. The resource GCN and task MLP are separate encoders.}
 \label{fig:overview}
\end{figure*}

\subsection{Initialization, Operators, and Restarts}
Initial solutions include a criticality-based hybrid deployment and an
earliest-finish-time deployment. When learned guidance is enabled, a GCN-led
greedy construction is added. Candidates are ranked by estimated feasibility,
violation, and penalized objective.

The destroy pool contains random removal, critical-task removal, high-risk
removal, congested-node removal, and removal of excessive backups. A learned
operator additionally targets tasks placed on low-scoring nodes. Five repair
operators emphasize normalized insertion cost, recovery delay, monetary cost,
active-copy enhancement, or passive-mode conversion. Removing and repairing
tasks changes both deployment locations and the conditions under which a
task receives replicas.

Destroy and repair operators are sampled by weighted roulette. Segment-level
updates combine their previous weights with rewards for accepted, improving,
or best-search outcomes. A candidate with a worse current search score can
still be accepted with simulated-annealing probability
\begin{equation}
 P_{\mathrm{acc}}=\exp\!\left(-\frac{S(\chi')-S(\chi)}{T}\right),
 \label{eq:sa}
\end{equation}
where $S$ is the active feasibility-first or penalized score, and the
temperature decreases geometrically. At prolonged stagnation, the solver
destroys and repairs a larger portion of a reference deployment, resets the
temperature, and consumes one of a bounded number of restarts. If the restart
budget is exhausted, it can stop before the iteration limit.

% IEEEtran uses a figure float for pseudocode instead of a dedicated
% algorithm float. algpseudocode supplies only the pseudocode layout.
\begin{figure}[!t]
\centering
\begin{minipage}{\columnwidth}
\hrule\vspace{3pt}
\footnotesize
\begin{algorithmic}[1]
\Require Workflow/resource instance, parameters, iteration limit $B$
\State Compute structural features and cached learned scores
\State Construct diverse initial solutions; select current $\chi$
\State Initialize operator weights, temperature, and restart budget
\State $\chi_F\gets\varnothing$; $J_F\gets+\infty$
\If{the quick screen accepts the initial solution}
  \State Fully verify it; set $\chi_F\gets\chi$, $J_F\gets J(\chi)$ if feasible
\EndIf
\For{$t=1,\ldots,B$}
  \State Select destroy and repair operators by adaptive roulette
  \State $\chi'\gets\operatorname{Repair}(\operatorname{Destroy}(\chi))$
  \State Evaluate $J_0$, nominal timing, $\rhat$, and $V$ for $\chi'$
  \State Use $V+\epsilon J_0$ until $\chi_F$ exists, then use $\Phi$
  \State Accept $\chi'$ into the current search state by improvement or SA
  \If{quick-feasible$(\chi')$ and $J_0(\chi')<J_F$}
    \State Enumerate all $\omega\in\Omega(\chi')$ with the shared verifier
    \If{$\chi'$ passes the full capacity and recovery gate}
      \State $\chi_F\gets\chi'$; $J_F\gets J(\chi')$
    \EndIf
  \EndIf
  \State Accumulate operator rewards; update weights at segment boundaries
  \If{stagnation threshold is reached}
    \State If permitted, restart with a large destroy/repair step; otherwise stop
    \State Evaluate the restart and apply the same feasible-register check
  \EndIf
\EndFor
\State \Return $\chi_F$ if present; otherwise return the best quick-search state
\end{algorithmic}
\vspace{3pt}\hrule
\end{minipage}
\caption{GCN-ALNS search procedure with a shared recovery verifier.}
\label{alg:gcnalns}
\end{figure}

\subsection{Resource-Graph Scoring Network}
The physical resource graph connects each node to its $k$ nearest neighbors
under unit-data communication delay, then symmetrizes the edges. With
$\widetilde A=A+I$ and its degree matrix $\widetilde D$, propagation uses
\begin{equation}
 H^{(\ell+1)}=\widetilde D^{-1/2}\widetilde A\widetilde D^{-1/2}
                  \bigl(H^{(\ell)}W^{(\ell)}+b^{(\ell)}\bigr),
 \label{eq:gcn}
\end{equation}
with ReLU and dropout between graph layers. Node inputs combine 12 resource
features with a learned fault-domain embedding. A separate MLP maps 12 task
features---including workload, rank, slack, degree, level, and structural
criticality---to a task embedding. No graph convolution is performed on the
workflow DAG itself.

Scalar criticality combines normalized upward rank, inverse structural slack,
and out-degree:
\begin{equation}
 h_i=\vartheta_1\widetilde{\mathrm{rank}}_i+
 \vartheta_2\widetilde{(\mathrm{slack}_i+\varepsilon)^{-1}}+
 \vartheta_3\widetilde{\mathrm{outdeg}}_i,
\end{equation}
with weights $(0.4,0.4,0.2)$ in the registered configuration. Structural
features are computed from average execution and communication estimates
before choosing a deployment.

For task $i$ and node $k$, a nine-dimensional interaction vector $z_{ik}$
contains execution time, a node-risk feature, a local recovery-risk proxy,
deployment cost, communication affinity, a compute-speed ratio, and three
resource/deadline compatibility indicators. The placement preference is
\begin{equation}
 s_{ik}=\operatorname{sigmoid}\!\left(
 \MLP([h_i^{T},h_k^{N},h_i^{T}\odot h_k^{N},z_{ik}])\right).
 \label{eq:placement_score}
\end{equation}
Here $h_i^{T}$ is a task embedding and must not be confused with scalar
criticality $h_i$. Predictions for all task--node
pairs are computed once and cached; they do not change with the partial
deployment during an ALNS run.

\subsection{Integrating the Learned Preferences}
For a task's current candidate node set, each insertion component is
min--max normalized within that set. Insertion minimizes
\begin{equation}
\begin{split}
 I_{ik}={}&\alpha_1\widetilde{\mathrm{ECT}}_{ik}
 +\alpha_2\widetilde{\mathrm{Cost}}_{ik}
 +\alpha_3\widetilde{\mathrm{Risk}}_{ik}\\
 &+\alpha_4\widetilde{\mathrm{RecRisk}}_{ik}
 +\alpha_5\widetilde{\mathrm{Load}}_{ik}
 -\alpha_6 h_i\widetilde s_{ik}.
\end{split}
\label{eq:insertion}
\end{equation}
Normalization also applies to the learned score. Candidate screening checks
resource capacity, relevant domain exclusions, and an estimated task finish
against the deadline. It does not establish a hard recovery-delay bound. A
fallback can still construct a candidate when screening finds no node, which
is another reason that final verification is necessary.

The learned scorer influences the greedy initial construction and the
low-score destroy operator as well as~\eqref{eq:insertion}. Replica selection
can also use its node preferences. Repair chooses active replication when a
local recovery estimate exceeds the effective budget. The final method
combines this heuristic protection rule with a learned placement scorer.
Disabling the learned
replica-ranking hook in the component study removes pair-score guidance for
backup destinations while retaining the heuristic protection rule.

\subsection{Learning from Feasible Demonstrations}
A teacher panel runs passive and replicated heuristics together with ALNS on
each designated learning or validation input. Every demonstration is checked
by the same scenario verifier. Among feasible panel outputs, the cheapest
top fraction $\zeta=0.2$ is retained, with at least one output if any is
feasible. Chosen task--node pairs become positive examples. A positive pair
from a demonstration is weighted by
\begin{equation}
 w_{ik}=(1+h_i)\left[1-\frac{\rmax}{\reff}\right]_{[0,1]},
 \label{eq:teacher_weight}
\end{equation}
with the zero-budget case handled separately. Repeated positives retain the
largest associated weight. Unselected compatible nodes provide sampled
contrastive negatives; they are not proven to be invalid assignments in all
possible deployments.

Training combines positive and negative binary cross-entropy with a pairwise
ranking term:
\begin{equation}
\begin{split}
 \mathcal L={}&\mathcal L_{\mathrm{BCE}}^{+}
 +\mathcal L_{\mathrm{BCE}}^{-}
 -\lambda_{\mathrm{rank}}
   \mathbb E_{(i,k^+,k^-)}
   [\overline w_{ik^+}\log\operatorname{sigmoid}(a_{ik^+}-a_{ik^-})],
\end{split}
\label{eq:training_loss}
\end{equation}
where $a_{ik}$ is a placement logit and positive weights are normalized to
mean one within a training batch. The final scorer is selected on the original
validation instances and frozen before held-out testing. Alternative training
configurations are retained as internal development records rather than
presented as separate solver variants.

\subsection{Verification Cost and Practical Limits}
If a deployment has $F$ enumerated faults, verification performs $F$ forward
replays. Each replay can scan the entire node set when selecting recovery
destinations and must evaluate the affected dependencies. The current Python
implementation also performs repeated edge scans when obtaining predecessor
lists; these costs should not be hidden behind a claim of constant-time
dependency lookup. Caching learned scores
avoids repeating a network forward pass at every insertion, but insertion,
schedule evaluation, and full verification can still dominate elapsed time.
Measured runtime and iteration-based search effort are therefore reported
separately.
